\label{sec:conclusion}

This paper presented an end-to-end methodology for wireless channel dataset
generation, fingerprint localization, and channel charting on the Otaniemi
campus.
A cropped $132\,\text{m}\times 147\,\text{m}$ sub-region of the Aalto
School of Electrical Engineering was instrumented with nine rooftop
5G NR~Band~n78 sites ($f_c = 3.6\,\text{GHz}$,
$4\times4$ antennas per BS), and a $2.5\,\text{m}$ UE grid with
3\,069 valid positions.
The wireless channel dataset was generated using MATLAB's Ray Tracing
toolbox with configurable path-tracing parameters (reflections, diffractions, frequency)
to produce physically realistic multipath channel characteristics for both
localization and channel-charting applications.

The resulting fingerprint matrices were used in two complementary ways.
A $3\,069 \times 63$ MATLAB-derived fingerprint matrix (with AoA, RSS, delay, and LOS features)
generated from MATLAB ray-tracing and processed in Python served as an independent
validation dataset.
For \emph{fingerprint localization} (Section~\ref{sec:localization}),
six estimators were benchmarked on the MATLAB-derived dataset under a 70\,\%/30\,\% train/test split;
wKNN with coarse-to-fine refinement and MLP regression reached comparable performance
($\text{MAE} \approx 18\,\text{m}$), demonstrating the robustness of fingerprinting
across different ray-tracing implementations and feature representations.
A classification head on the same grid failed because the cell count
exceeded the number of training samples --- a clear demonstration of
why regression is preferable to classification for dense grids.
For \emph{channel charting} (Section~\ref{sec:channel_charting}), four
DR methods (PCA, t-SNE, autoencoder, UMAP) were scored with
TW/CT/KS/NNLE on both the Otaniemi\_small and MATLAB-derived datasets;
t-SNE produced strong unsupervised maps on both datasets, with
$\text{TW}=0.918$, $\text{CT}=0.958$ on Otaniemi\_small and $\text{TW}=0.975$, $\text{CT}=0.945$
on the MATLAB-derived dataset --- comparable to the best supervised regressor without
using any UE labels.

Overall the Otaniemi campus testbed, validated through independent MATLAB ray-tracing
simulation, is a useful mid-scale environment: large enough to produce
non-trivial NLOS structure that stresses both the localizer and the chart, yet tractable
for CPU-only processing.
The main limitations observed were (i) the dominant geometric ambiguity
in the shadow regions behind large facades, which inflates the tail of
the localization error CDF across both datasets; (ii) the reduced feature
dimensionality in the MATLAB-derived feature vector (63 dimensions: AoA, RSS, delay, LOS)
compared to full CSI fingerprints, which limits performance for data-hungry models;
and (iii) the limited number of grid samples for deep learning models such
as the plain MSE autoencoder.

\subsection*{Future Work}

Several follow-on directions are natural.
First, the fingerprint vector currently uses only amplitude-domain
features; incorporating the per-path AoA/AoD angles already available in
the HDF5 dataset is expected to sharpen both the localizer and the
chart~\cite{hoydis2023sionna_rt,3gpp_nr_channel}.
Second, replacing the plain-MSE autoencoder with a
Siamese/triplet-loss AE~\cite{lei2019siamese_cc} would give the latent
space an explicit metric-preservation objective and should close most of
the gap to t-SNE.
Third, extending the scene size --- by reverting to the full 77-building
Otaniemi model once a GPU host is available --- would capture more
propagation diversity and enable multi-point channel charting across
adjacent cells.
Fourth, a semi-supervised channel-chart formulation, in which a small
subset of labelled anchors regularises the latent space, is a direct
avenue for the extra-credit item in the course
assignment~\cite{studer2018channel_charting}.
Finally, closing the loop by using the chart as an input to a
beam-management or handover policy would demonstrate the practical
utility of the learned radio-space embedding for real 5G/beyond-5G
systems~\cite{sionna2022}.
